% ======================================================================
\section{Results}
\label{sec:results}
Every number in this section is generated from the released result files; the prose quotes
macros, never typed values.

\subsection{Accuracy on eight national systems}
\label{sec:grids}
\tabv{tab_grids_v2}
\tabv{tab_grids_v2_b}
\begin{figure}[!t]
\centering
\figv{fig_horizon_v2}{\columnwidth}
\caption{Mean absolute error by lead time of the two-week ($H=336$) models under regime R2,
per system, in MW (seed means of the per-step test error of the headline methods).}
\label{fig:horizon}
\end{figure}
\begin{figure*}[!t]
\centering
\figv{fig_mcs_v2}{\textwidth}
\caption{Model Confidence Set membership ($\alpha=\nv{mcs.alpha}$, stationary bootstrap with
$B=\nv{mcs.B}$ replicates, block length at least the horizon) for every system--horizon cell
under R2. A filled cell means the method cannot be distinguished from the best in that cell;
the number is the method's WAPE rank in that cell. Rows are ordered by the number of cells
retained.}
\label{fig:mcs}
\end{figure*}

Table~\ref{tab:grids-v2} reports the WAPE of every learned forecaster on every system and
horizon under the headline regime R2: \nv{grid.n_models} methods on \nv{grid.n_cells}
system--horizon cells, each a three-seed mean. Figure~\ref{fig:horizon} plots the headline
methods against horizon and Fig.~\ref{fig:mcs} shows, cell by cell, which methods the Model
Confidence Set retains. Four features of the table survive any reasonable reading of it.

\paragraph{One method leads} TiDE, the covariate-native MLP, is the single most accurate
method in \nv{grid.TiDE.n_best} of the \nv{grid.n_cells} cells, has a mean WAPE rank of
\nv{grid.TiDE.meanrank_wape}, and is retained by the confidence set in
\nv{mcs.grid.TiDE.n_in} cells. We state this first because the paper proposes a different
model, and the reader should know at the outset that the proposed model is not the most
accurate point forecaster in this benchmark. Section~\ref{sec:info} and the discussion locate
the difference: the two models receive the same information and differ in the pathway by which
the weather enters.

\paragraph{The proposed model is second, and the top is a set} PLANET has the next mean rank,
\nv{grid.PLANET.meanrank_wape}, and is retained by the confidence set in \nv{mcs.grid.PLANET.n_in} cells
(\nv{mcs.grid.PLANET.H24.n_in} at one day, \nv{mcs.grid.PLANET.H96.n_in} at four days and
\nv{mcs.grid.PLANET.H336.n_in} at two weeks, where the sets retain between
\nv{mcs.grid.H336.size_min} and \nv{mcs.grid.H336.size_max} of the \nv{grid.n_models} methods
against \nv{mcs.grid.H24.size_min} to \nv{mcs.grid.H24.size_max} at one day); it is followed
by PatchTST+Cal
(\nv{grid.CalTST.meanrank_wape}, retained in \nv{mcs.grid.CalTST.n_in} cells), PLANET-X
(\nv{grid.PLANETX.meanrank_wape}, \nv{mcs.grid.PLANETX.n_in}), N-BEATS
(\nv{grid.NBEATS.meanrank_wape}, \nv{mcs.grid.NBEATS.n_in}), TimeMixer
(\nv{grid.TimeMixer.meanrank_wape}, \nv{mcs.grid.TimeMixer.n_in}), N-HiTS
(\nv{grid.NHITS.meanrank_wape}, \nv{mcs.grid.NHITS.n_in}) and PatchTST
(\nv{grid.PatchTST.meanrank_wape}, \nv{mcs.grid.PatchTST.n_in}). Pairwise, PLANET is
significantly more accurate than the compared method in \nv{dm.total.wins} of
\nv{dm.n_tests} Holm-adjusted Diebold--Mariano tests, indistinguishable in \nv{dm.total.ties}
and significantly less accurate in \nv{dm.total.losses} (with the stationary-bootstrap
variance of Section~\ref{sec:stats}: \nv{dm.total.wins_boot}, \nv{dm.total.ties_boot} and
\nv{dm.total.losses_boot}; \nv{dm.total.wins_allseeds} of the wins and
\nv{dm.total.losses_allseeds} of the losses hold in each of the three seeds taken separately);
\nv{dm.TiDE.losses} of those losses
are to TiDE, and no other method is significantly better than PLANET in more than
\nv{dm.CalTST.losses} cells. Against the compact models nearest to it the record
(wins/ties/losses) is \nv{dm.CalTST.wins}/\nv{dm.CalTST.ties}/\nv{dm.CalTST.losses} against
PatchTST+Cal, \nv{dm.NBEATS.wins}/\nv{dm.NBEATS.ties}/\nv{dm.NBEATS.losses} against N-BEATS
and \nv{dm.PatchTST.wins}/\nv{dm.PatchTST.ties}/\nv{dm.PatchTST.losses} against PatchTST;
against TiDE it is \nv{dm.TiDE.wins}/\nv{dm.TiDE.ties}/\nv{dm.TiDE.losses}. The claim the table
supports is therefore not that the tier dominates, but that a model built around it is, after
TiDE, the most consistently accurate of \nv{grid.n_models} methods, ahead of every other
Transformer, linear, basis-expansion and library model under the same inputs and the same
budget.

\paragraph{Where the margin lies} Relative to the strongest \emph{other} method in each cell,
whichever it is, PLANET's WAPE is \nv{ref.grid.mean_gain_wape_pct} on average (positive means
PLANET is better), between \nv{ref.grid.min_gain_wape_pct} and
\nv{ref.grid.max_gain_wape_pct} across cells; by horizon the figure is
\nv{ref.grid.H24.mean_gain_wape_pct} at one day, \nv{ref.grid.H96.mean_gain_wape_pct} at four
days and \nv{ref.grid.H336.mean_gain_wape_pct} at two weeks, so the deficit to TiDE opens
beyond one day rather than closing. Against TiDE directly, PLANET is ahead on Germany at all
three horizons (\nv{ref.OPSD-DE.H24.vs.TiDE.gain_wape_pct},
\nv{ref.OPSD-DE.H96.vs.TiDE.gain_wape_pct} and \nv{ref.OPSD-DE.H336.vs.TiDE.gain_wape_pct}),
within three percent of it on France at one day (\nv{ref.OPSD-FR.H24.vs.TiDE.gain_wape_pct})
and on Great Britain at every horizon (\nv{ref.OPSD-GB.H24.vs.TiDE.gain_wape_pct},
\nv{ref.OPSD-GB.H96.vs.TiDE.gain_wape_pct}, \nv{ref.OPSD-GB.H336.vs.TiDE.gain_wape_pct}), and
behind elsewhere, by the most on Poland at four days
(\nv{ref.OPSD-PL.H96.vs.TiDE.gain_wape_pct}) and France at two weeks
(\nv{ref.OPSD-FR.H336.vs.TiDE.gain_wape_pct}). Section~\ref{sec:discussion} relates this
pattern to how much of each system's load the weather explains. Peak-height error, the
operator's costliest error, orders the methods the same way: averaged over the cells the peak
MAPE is \nv{grid.TiDE.mean_peak} for TiDE, \nv{grid.PLANET.mean_peak} for PLANET,
\nv{grid.CalTST.mean_peak} for PatchTST+Cal, \nv{grid.NBEATS.mean_peak} for N-BEATS and
\nv{grid.PatchTST.mean_peak} for PatchTST, with mean peak-timing errors of
\nv{grid.TiDE.mean_dt}, \nv{grid.PLANET.mean_dt}, \nv{grid.CalTST.mean_dt},
\nv{grid.NBEATS.mean_dt} and \nv{grid.PatchTST.mean_dt} hours.

\paragraph{Seed stability} The mean seed standard deviation of WAPE over the cells is
\nv{grid.PLANET.mean_wape_sd} for PLANET, \nv{grid.CalTST.mean_wape_sd} for PatchTST+Cal,
\nv{grid.NBEATS.mean_wape_sd} for N-BEATS, \nv{grid.TiDE.mean_wape_sd} for TiDE,
\nv{grid.PatchTST.mean_wape_sd} for PatchTST and \nv{grid.TFT.mean_wape_sd} for TFT: the compact models are all stable, PLANET marginally the most so, and TFT the least of the
six. For an operator who trains one model and deploys it, the
seed spread is part of the expected error, and the table reports it beside every mean for that
reason.

\paragraph{The rest of the field} The long-horizon Transformers of the first generation are
not competitive on national load under the full-information regime, the only one they were
run under: Informer and Autoformer have mean WAPE
ranks of \nv{grid.Informer.meanrank_wape} and \nv{grid.Autoformer.meanrank_wape} of
\nv{grid.n_models}, and Informer is on average behind the seasonal-na\"ive reference
(\nv{grid.SeasonalNaive.meanrank_wape}). The one-layer linear family sits in the middle of the
field (NLinear \nv{grid.NLinear.meanrank_wape}, RLinear \nv{grid.RLinear.meanrank_wape}, FITS
\nv{grid.FITS.meanrank_wape}); DLinear, whose moving-average decomposition is mis-specified
for a signal whose ``trend'' is the weekly cycle, is the weakest of them
(\nv{grid.DLinear.meanrank_wape}). Among the covariate-native library models TiDE
(\nv{grid.TiDE.meanrank_wape}) and N-HiTS (\nv{grid.NHITS.meanrank_wape}) are the strongest and
TFT (\nv{grid.TFT.meanrank_wape}) and TimeXer (\nv{grid.TimeXer.meanrank_wape}) the weakest,
under the same tuning budget as every other model. Among the classical methods, Hong's regression with an autoregressive model of its
residual ranks highest (mean rank \nv{grid.TaoRefitAR.meanrank_wape}; grid-mean WAPE
\nv{grid.TaoRefitAR.mean_wape} against \nv{grid.PLANET.mean_wape} for PLANET), the
gradient-boosted trees, given the same information as the shared head, come next
(\nv{grid.GBRT.meanrank_wape}), and the same regression fitted once or monthly without the
residual model is the weakest learned method in the study (grid-mean WAPE
\nv{grid.TaoRefit.mean_wape}), below the seasonal naive on several systems because its trend
term extrapolates a level that the test span does not follow (Section~\ref{sec:setup}). The
regression with the residual model is therefore the classical benchmark the deep models have
to beat; it remains behind the leading deep models in every cell, which is the case for the deep methods that
Section~\ref{sec:info} makes on information-controlled terms.

\subsection{Value of information}
\label{sec:info}
\tabv{tab_regimes_v2}
\tabv{tab_regimes_v2_b}
\begin{figure*}[!t]
\centering
\figv{fig_regimes_v2}{\textwidth}
\caption{WAPE of each headline method under the three information regimes, one panel per
horizon, mean over systems (dots: individual systems). PLANET and PatchTST+Cal carry the
calendar in their design matrix, so for them R0 is not a past-only regime and the R0--R1 step
measures only what the shared head adds beyond the built-in basis.}
\label{fig:regimes}
\end{figure*}

Table~\ref{tab:regimes-v2} and Fig.~\ref{fig:regimes} report the headline methods under the
three regimes. Two readings matter. The first is the size of the effect: over the headline
methods without a built-in calendar, the calendar block reduces WAPE by
\nv{info.grid.R1_vs_R0.mean_dwape_pct} on average and the weather block by a further
\nv{info.grid.R2_vs_R1.mean_dwape_pct}, for a combined \nv{info.grid.R2_vs_R0.mean_dwape_pct}
from R0 to R2 (N-BEATS \nv{info.grid.NBEATS.R2_vs_R0.mean_dwape_pct}, PatchTST
\nv{info.grid.PatchTST.R2_vs_R0.mean_dwape_pct}, DLinear
\nv{info.grid.DLinear.R2_vs_R0.mean_dwape_pct}, TiDE \nv{info.grid.TiDE.R2_vs_R0.mean_dwape_pct},
N-HiTS \nv{info.grid.NHITS.R2_vs_R0.mean_dwape_pct}, TFT
\nv{info.grid.TFT.R2_vs_R0.mean_dwape_pct}, TimeXer \nv{info.grid.TimeXer.R2_vs_R0.mean_dwape_pct},
gradient-boosted trees \nv{info.grid.GBRT.R2_vs_R0.mean_dwape_pct}). That change, obtained by
handing a fixed architecture the inputs an operator already has, is larger than the gap between
PLANET and the median learned method within R2 (\nv{grid.median_gap_wape_pct}, PLANET relative
to the median). A benchmark that lets each method consume whatever its interface accepts is
therefore mostly measuring interfaces.

The second reading concerns the two calendar-aware models and TiDE. PLANET's basis already
encodes local time, the weekly cycle and the holiday day-types, so under R0 it is not a
past-only forecaster: on Germany at four days its R0 error is \nv{R0.OPSD-DE.PLANET.H96.wape}
against \nv{R0.OPSD-DE.NBEATS.H96.wape} for N-BEATS, \nv{R0.OPSD-DE.PatchTST.H96.wape} for
PatchTST and \nv{R0.OPSD-DE.TiDE.H96.wape} for TiDE, which measures the value of the built-in
calendar itself, and PatchTST+Cal, which learns the same design matrix globally, reaches
\nv{R0.OPSD-DE.CalTST.H96.wape}. The R0--R1 step then adds the calendar block through the shared head: averaged over the
grid, \nv{info.grid.PLANET.R1_vs_R0.mean_dwape_pct} for PLANET and
\nv{info.grid.CalTST.R1_vs_R0.mean_dwape_pct} for PatchTST+Cal, so a globally learned calendar
correction still adds to the per-window fit, while a model that already learns the calendar
map globally gains nothing from a second copy. The weather step is where TiDE separates: it is
\nv{info.grid.PLANET.R2_vs_R1.mean_dwape_pct} for PLANET, close to the
\nv{info.grid.R2_vs_R1.mean_dwape_pct} mean of the models that receive the block through the
head, and \nv{info.grid.TiDE.R2_vs_R1.mean_dwape_pct} for TiDE, which reads it natively; on
France at four days the two figures are \nv{info.OPSD-FR.PLANET.H96.R2_vs_R1.dwape_pct} and
\nv{info.OPSD-FR.TiDE.H96.R2_vs_R1.dwape_pct}, on Germany
\nv{info.OPSD-DE.PLANET.H96.R2_vs_R1.dwape_pct} and
\nv{info.OPSD-DE.TiDE.H96.R2_vs_R1.dwape_pct}. The tier does not compete with the covariates,
it models what they do not explain; the covariate pathway decides how much is left to explain,
and that is the gap of Table~\ref{tab:grids-v2}. The weather block is not a gain for every model on the calendar-driven systems: DLinear is
worse with it on Germany at every horizon (\nv{info.OPSD-DE.DLinear.H24.R2_vs_R1.dwape_pct},
\nv{info.OPSD-DE.DLinear.H96.R2_vs_R1.dwape_pct} and
\nv{info.OPSD-DE.DLinear.H336.R2_vs_R1.dwape_pct}) and N-HiTS on the Netherlands, and one
model is harmed outright. Hong's regression is fitted on observed temperatures and forecasts with the degraded ones, as an
operator uses it, whereas the neural models are trained on the degraded block and learn to
discount it; on Germany, where the temperature carries little load, the forecast error that
its cubic temperature terms amplify exceeds the signal they add, and the regression is more
accurate without them (\nv{R1.OPSD-DE.TaoVanilla.H24.wape} against
\nv{OPSD-DE.TaoVanilla.H24.wape} WAPE at one day fitted once,
\nv{R1.OPSD-DE.TaoRefitAR.H24.wape} against \nv{OPSD-DE.TaoRefitAR.H24.wape} with the
residual model), while on France the same block nearly halves the error of the once-fitted
model (\nv{R1.OPSD-FR.TaoVanilla.H24.wape} against \nv{OPSD-FR.TaoVanilla.H24.wape}).

\subsection{Day-ahead forecasts against the operators}
\label{sec:dayahead}
\tabv{tab_dayahead_v2}
\tabv{tab_dayahead_v2_b}
\tabv{tab_dayahead_v2_c}
\begin{figure}[!t]
\centering
\figv{fig_dayahead_v2}{\columnwidth}
\caption{Day-ahead WAPE of each learned model (seed mean of the single-seed models at the
protocol lookback) and of the transmission operator's own published forecast (bar) on the same
days and target hours, per system, systems sorted by the operator's error. Origin 11:00 UTC on
day $D-1$, scored 00:00--23:00 UTC of day $D$.}
\label{fig:dayahead}
\end{figure}

Table~\ref{tab:dayahead-v2} holds every learned model to the operator's information cut-off
and scores the operator's published forecast on the same days and hours. The published series
of \nv{da.grid.biased_systems} sit \nv{da.OPSD-GB.TSO.bias_pct} and
\nv{da.OPSD-NL.TSO.bias_pct} below the published actual load on the scored days
(Section~\ref{sec:dayahead-proto}), so the comparison is read on the \nv{da.grid.n_unbiased}
systems whose operator bias is within \nv{da.grid.bias_tol_pct}. On those PLANET is more
accurate than the operator on \nv{da.grid.PLANET.n_beats_tso_unbiased}: Germany, where the
operator's WAPE over the \nv{da.OPSD-DE.n_days} scored days is \nv{tso.OPSD-DE.wape} against
\nv{da.OPSD-DE.PLANET.wape} for PLANET, a change of \nv{da.OPSD-DE.PLANET.gain_wape_pct} with
the peak-height error changed by \nv{da.OPSD-DE.PLANET.gain_peak_pct}; Poland
(\nv{da.OPSD-PL.PLANET.gain_wape_pct}); and Belgium (\nv{da.OPSD-BE.PLANET.gain_wape_pct}). The
operator is ahead on France (\nv{da.OPSD-FR.PLANET.gain_wape_pct}), Italy
(\nv{da.OPSD-IT.PLANET.gain_wape_pct}) and, by a wide margin, Spain
(\nv{da.OPSD-ES.PLANET.gain_wape_pct}), whose published forecast is the most accurate in the
corpus (\nv{tso.OPSD-ES.wape} WAPE). TiDE beats the operator on
\nv{da.grid.TiDE.n_beats_tso_unbiased} of the six, adding Italy
(\nv{da.OPSD-IT.TiDE.gain_wape_pct}); N-BEATS and PatchTST+Cal on
\nv{da.grid.NBEATS.n_beats_tso_unbiased} and \nv{da.grid.CalTST.n_beats_tso_unbiased}; and on \nv{da.grid.any.n_beats_tso_unbiased} of the six at least one learned model or seed
ensemble does (single models reach \nv{da.grid.TiDE.n_beats_tso_unbiased} at most). The
equal-weight mean of the PLANET and TiDE seed ensembles beats the operator on
\nv{da.grid.Comb2-ens.n_beats_tso_unbiased} of the six (significantly on
\nv{da.grid.Comb2-ens.n_sig_beats_tso_unbiased}), against
\nv{da.grid.PLANET-ens.n_beats_tso_unbiased} for the PLANET ensemble and
\nv{da.grid.TiDE-ens.n_beats_tso_unbiased} for TiDE's (\nv{da.grid.PLANET-ens.n_sig_beats_tso_unbiased}
and \nv{da.grid.TiDE-ens.n_sig_beats_tso_unbiased} significantly; combination rows of
Table~\ref{tab:dayahead-v2}). Refitting PLANET and TiDE from scratch at the start of every
test month (refit rows) changes their WAPE on the six by
\nv{da.grid.PLANET-refit.mean_dwape_vs_once_unbiased_pct} and
\nv{da.grid.TiDE-refit.mean_dwape_vs_once_unbiased_pct} on average against the once-trained run
of the same seed and leaves them ahead of the operator on
\nv{da.grid.PLANET-refit.n_beats_tso_unbiased} and \nv{da.grid.TiDE-refit.n_beats_tso_unbiased}
systems, so training once is not what separates the learned models from the operators. On the
holiday days of the test span the PLANET and TiDE ensembles are ahead of the operator on
\nv{da.grid.PLANET-ens.n_beats_tso_hol_unbiased} and \nv{da.grid.TiDE-ens.n_beats_tso_hol_unbiased}
of the six systems, inside the first-wave lockdown window on
\nv{da.grid.PLANET-ens.n_beats_tso_lock_unbiased} and \nv{da.grid.TiDE-ens.n_beats_tso_lock_unbiased},
and outside it on \nv{da.grid.PLANET-ens.n_beats_tso_xlock_unbiased} and
\nv{da.grid.TiDE-ens.n_beats_tso_xlock_unbiased}: what the operators keep is concentrated on
the holiday days, where their mean WAPE over the six is \nv{da.grid.tso.mean_wape_hol_unbiased}
against \nv{da.grid.PLANET-ens.mean_wape_hol_unbiased} and
\nv{da.grid.TiDE-ens.mean_wape_hol_unbiased} for the two ensembles, against
\nv{da.grid.tso.mean_wape_xlock_unbiased}, \nv{da.grid.PLANET-ens.mean_wape_xlock_unbiased} and
\nv{da.grid.TiDE-ens.mean_wape_xlock_unbiased} on the ordinary days outside the lockdown
window. On the two
biased systems every learned model is far ahead of the published series, which says nothing
about the operators. The operators' forecasts are the product of a process we do not
reproduce, with access to the latest weather run and, in some systems, to manual adjustment
for known events, while the learned models hold degraded reanalysis rather than an archived
operational forecast. We read the result as evidence that the protocol is fair, and that the
accuracy an operator obtains is reachable by a model trained under it on half of the systems,
not as a claim to replace the operator.

\subsection{Component ablation}
\label{sec:ablation}
\tabv{tab_ablation_wape_v2}
\begin{figure}[!t]
\centering
\figv{fig_ablation_v2}{\columnwidth}
\caption{Relative change in test WAPE when each component of PLANET is removed, on both
residual backbones: mean over the six DE and FR system--horizon cells with a 95\,\%
confidence interval over cells; small markers are the individual cells. Positive means the
component helps.}
\label{fig:ablation}
\end{figure}

Table~\ref{tab:ablation-wape-v2} removes one component at a time on both residual backbones,
on the two reference systems, at three horizons and over three seeds (five for the full model
and the four principal variants on the patch backbone), each run paired with the full model of
the same seed on the same backbone (\nv{abl.n_runs} runs). The first row is the component the
paper proposes. Removing the phase-locked tier, which leaves the backbone alone with the
shared head, the instance normalisation and the climatological tier, raises WAPE by
\nv{abl.patch.no_pl.mean_dwape_pct} on the patch Transformer and by
\nv{abl.basis.no_pl.mean_dwape_pct} on the basis stack, averaged over the six cells. The average
hides a structure that the per-cell numbers show and that the rest of the paper confirms: on
Germany the tier is worth \nv{abl.patch.OPSD-DE.no_pl.H24.dwape_pct} and
\nv{abl.patch.OPSD-DE.no_pl.H96.dwape_pct} at one and four days on the patch backbone
(\nv{abl.basis.OPSD-DE.no_pl.H24.dwape_pct} and \nv{abl.basis.OPSD-DE.no_pl.H96.dwape_pct} on
the basis stack), on France \nv{abl.patch.OPSD-FR.no_pl.H24.dwape_pct} and
\nv{abl.patch.OPSD-FR.no_pl.H96.dwape_pct} (\nv{abl.basis.OPSD-FR.no_pl.H24.dwape_pct} and
\nv{abl.basis.OPSD-FR.no_pl.H96.dwape_pct}), and at two weeks it is worth
\nv{abl.patch.OPSD-DE.no_pl.H336.dwape_pct} and \nv{abl.patch.OPSD-FR.no_pl.H336.dwape_pct}
on the patch backbone. The two backbones share no code, no inductive bias and no
hyper-parameter, and the component moves both in the same direction in every cell but two, France at one
day on the basis stack (\nv{abl.basis.OPSD-FR.no_pl.H24.dwape_pct}) and at two weeks on the
patch backbone (\nv{abl.patch.OPSD-FR.no_pl.H336.dwape_pct}), where its effect on that
backbone is within a point of zero, so it is doing something neither does on its
own; what it does is worth most on the German system and, on the patch backbone, at the
horizons short enough for the per-window coefficients to hold, and
Section~\ref{sec:discussion} shows that the systems' weather dependence does not predict where
it pays. The
converse row is as informative: the closed form alone, without any residual network, costs
\nv{abl.patch.no_resid.mean_dwape_pct} and \nv{abl.basis.no_resid.mean_dwape_pct}, so the two
parts are complementary, and the residual network is the larger of the two.

The design choices of Section~\ref{sec:basis} are then tested against their nearest
alternatives, and here the ablation is less kind to the introduction than we expected. Locking
the basis to the window instead of the calendar costs \nv{abl.patch.relative.mean_dwape_pct}
and \nv{abl.basis.relative.mean_dwape_pct} on average, and
\nv{abl.patch.OPSD-DE.relative.H96.dwape_pct} and
\nv{abl.basis.OPSD-DE.relative.H96.dwape_pct} on Germany at four days, the cell where the tier
itself is worth most; on France the choice is immaterial. Evaluating the basis on UTC instead
of local civil time changes WAPE by \nv{abl.patch.utc.mean_dwape_pct} and
\nv{abl.basis.utc.mean_dwape_pct}: the one-hour phase error of half the year is repaired by
the residual network, and the local-time argument of the introduction, while correct, has no
measurable consequence at hourly resolution. Removing the day-type interactions costs
\nv{abl.patch.no_daytype.mean_dwape_pct} on the patch backbone and
\nv{abl.basis.no_daytype.mean_dwape_pct} on the basis stack; removing only the holiday
day-type costs \nv{abl.patch.daytype_wk.mean_dwape_pct} and
\nv{abl.basis.daytype_wk.mean_dwape_pct}. Conditioning the residual network on the solved
coefficients changes WAPE by \nv{abl.patch.no_beta_token.mean_dwape_pct} and
\nv{abl.basis.no_beta_token.mean_dwape_pct} when removed, and on the basis backbone the
feature-wise affine form of that conditioning is \nv{abl.basis.basis_film.mean_dwape_pct}
relative to the projected form.

Three optional components are decided here rather than assumed, and none of them earns its
place. Replacing the learned observation weights by uniform weights changes WAPE by
\nv{abl.patch.uniform_w.mean_dwape_pct} and \nv{abl.basis.uniform_w.mean_dwape_pct}, removing
the climatological tier by \nv{abl.patch.no_slow.mean_dwape_pct} and
\nv{abl.basis.no_slow.mean_dwape_pct}, and removing the local trend column by
\nv{abl.patch.no_trend.mean_dwape_pct} and \nv{abl.basis.no_trend.mean_dwape_pct}; the lean
composition without all three is within \nv{abl.patch.lean.mean_dwape_pct} and
\nv{abl.basis.lean.mean_dwape_pct} of the full model. The headline model is the full
composition, fixed before the grid was run; the ablation reports which parts could be removed
without loss, and the composition is not re-selected on the test set.

\subsection{When per-window estimation pays}
\label{sec:mech}
\begin{figure*}[!t]
\centering
\figv{fig_mech_v2}{\textwidth}
\caption{The three conditions under which a closed-form per-window fit should differ from a
learned calendar map, $H=96$. \textbf{(a)} Learning curves: test WAPE against the fraction of
the training split used (most recent windows), mean over DE, FR and ES. \textbf{(b)} Temporal
drift: monthly test WAPE of a model trained on the \emph{oldest} half of the split, against
the months elapsed since the end of its training data. \textbf{(c)} Zero-shot cross-grid
transfer: a model trained on one system (row) and evaluated, without fine-tuning, on another
(column); cells are WAPE, the diagonal is the in-domain model.}
\label{fig:mech}
\end{figure*}

The second contribution is a claim about \emph{when} the tier pays, made as three experiments
that could each have gone the other way (Fig.~\ref{fig:mech}). All three use the four-day
horizon, the tuned configuration of every model, three seeds, and seven of the headline
methods spanning the architectural families, N-BEATS, PatchTST, PatchTST+Cal, PLANET,
PLANET-X, TFT and TiDE, on Germany, France and Spain. The controlled comparison inside each is
PLANET against PatchTST+Cal, which receives the same design matrix and learns its coefficients
globally instead of solving them per window.

\paragraph{Scarce training data} With one tenth of the training split, about five months,
every method loses much of its accuracy: PLANET's WAPE rises by
\nv{mech.lc.grid.PLANET.f0.1.mean_dwape_pct} over its full-data value, PatchTST+Cal's by
\nv{mech.lc.grid.CalTST.f0.1.mean_dwape_pct}, PatchTST's by
\nv{mech.lc.grid.PatchTST.f0.1.mean_dwape_pct}, N-BEATS's by
\nv{mech.lc.grid.NBEATS.f0.1.mean_dwape_pct}, TFT's by
\nv{mech.lc.grid.TFT.f0.1.mean_dwape_pct}, TiDE's by \nv{mech.lc.grid.TiDE.f0.1.mean_dwape_pct}
and PLANET-X's by \nv{mech.lc.grid.PLANETX.f0.1.mean_dwape_pct}, averaged over the three
systems. At one quarter the figures are \nv{mech.lc.grid.PLANET.f0.25.mean_dwape_pct},
\nv{mech.lc.grid.CalTST.f0.25.mean_dwape_pct}, \nv{mech.lc.grid.PatchTST.f0.25.mean_dwape_pct},
\nv{mech.lc.grid.NBEATS.f0.25.mean_dwape_pct}, \nv{mech.lc.grid.TFT.f0.25.mean_dwape_pct},
\nv{mech.lc.grid.TiDE.f0.25.mean_dwape_pct} and \nv{mech.lc.grid.PLANETX.f0.25.mean_dwape_pct},
and at one half \nv{mech.lc.grid.PLANET.f0.5.mean_dwape_pct},
\nv{mech.lc.grid.CalTST.f0.5.mean_dwape_pct}, \nv{mech.lc.grid.PatchTST.f0.5.mean_dwape_pct},
\nv{mech.lc.grid.NBEATS.f0.5.mean_dwape_pct}, \nv{mech.lc.grid.TFT.f0.5.mean_dwape_pct},
\nv{mech.lc.grid.TiDE.f0.5.mean_dwape_pct} and \nv{mech.lc.grid.PLANETX.f0.5.mean_dwape_pct}.
PLANET degrades least of the seven from one quarter upwards; at one tenth the two lightest
architectures, PatchTST and N-BEATS, degrade least, and TiDE and PLANET-X most. In absolute
terms the picture is the one the headline table would predict: on Germany and Spain PLANET
has the lowest WAPE of the seven at one quarter (\nv{mech.lc.OPSD-DE.PLANET.H96.f0.25.wape}
against \nv{mech.lc.OPSD-DE.TiDE.H96.f0.25.wape} for TiDE on Germany,
\nv{mech.lc.OPSD-ES.PLANET.H96.f0.25.wape} against \nv{mech.lc.OPSD-ES.TiDE.H96.f0.25.wape} on
Spain), while on France TiDE keeps its lead at every fraction
(\nv{mech.lc.OPSD-FR.TiDE.H96.f0.25.wape} against \nv{mech.lc.OPSD-FR.PLANET.H96.f0.25.wape}
at one quarter). A calendar map has to be learned from examples of every day-type at every
hour, and a fortnight of holidays appears once per year of training data; a basis evaluated
at the calendar position needs no example of a holiday to know that one is coming, and only
the few coefficients of Eq.~\eqref{eq:solution} to know how the system responds to it. The advantage over the learned-calendar control is present on all three systems at both
fractions and clearest with half the history on Germany
(\nv{mech.lc.OPSD-DE.PLANET.H96.f0.25.wape} against \nv{mech.lc.OPSD-DE.CalTST.H96.f0.25.wape}
WAPE at a quarter, \nv{mech.lc.OPSD-DE.PLANET.H96.f0.5.wape} against
\nv{mech.lc.OPSD-DE.CalTST.H96.f0.5.wape} at a half); in absolute terms PLANET, which leads TiDE
on Germany with the full history, stays ahead of it there, passes it on Spain with a quarter
of the history, and does not on France, whose weather response no calendar basis supplies.

\paragraph{Transfer to an unseen system} A model trained on one national system and applied,
without any fine-tuning, to another sees a different level, a different daily shape, a
different holiday calendar and a different weather response. The target system's training
mean and scale are used for normalisation, as an operator deploying the model would have
them; the per-window solve absorbs what normalisation leaves of the level and the daily shape
by construction, the holiday day-types absorb the calendar, and only the weather response must
be learned. Averaged over the six ordered pairs of DE, FR and ES, the WAPE of the transferred
PLANET is \nv{mech.xfer.grid.PLANET.mean_dwape_pct} above the value its in-domain counterpart
reaches on the target, against \nv{mech.xfer.grid.TiDE.mean_dwape_pct} for TiDE,
\nv{mech.xfer.grid.CalTST.mean_dwape_pct} for PatchTST+Cal,
\nv{mech.xfer.grid.PatchTST.mean_dwape_pct} for PatchTST,
\nv{mech.xfer.grid.TFT.mean_dwape_pct} for TFT and \nv{mech.xfer.grid.NBEATS.mean_dwape_pct}
for N-BEATS; PLANET-X, which is not better than PLANET in domain, transfers best of all at
\nv{mech.xfer.grid.PLANETX.mean_dwape_pct}, a result we report as observed. PLANET and TiDE
lose the least of the established methods; PatchTST+Cal, which learned the same design matrix
on the source, loses more; and the models with no calendar representation beyond the head lose
most. The controlled pair, PLANET against PatchTST+Cal, is the mechanism claim in its cleanest
form, and TiDE's result says that a native covariate pathway transfers about as well as a solved
calendar does in relative terms (\nv{mech.xfer.grid.TiDE.mean_dwape_pct} against
\nv{mech.xfer.grid.PLANET.mean_dwape_pct}) and to a lower absolute error
(\nv{mech.xfer.grid.TiDE.mean_wape} against \nv{mech.xfer.grid.PLANET.mean_wape} WAPE). The
transferred PLANET thus reaches \nv{mech.xfer.grid.PLANET.mean_wape} WAPE at four days without seeing a single target window in
training, which is better than the in-domain value of Informer and of Hong's regression
without the residual model on the same systems in Table~\ref{tab:grids-v2}.

\paragraph{Temporal drift without retraining} The third experiment trains on the oldest
fraction of the split and evaluates on the same test span, so the model is, in effect, stale
by up to four years. With the oldest tenth PLANET's WAPE rises by
\nv{mech.drift.grid.PLANET.f0.1.mean_dwape_pct}, PatchTST+Cal's by
\nv{mech.drift.grid.CalTST.f0.1.mean_dwape_pct}, PatchTST's by
\nv{mech.drift.grid.PatchTST.f0.1.mean_dwape_pct}, N-BEATS's by
\nv{mech.drift.grid.NBEATS.f0.1.mean_dwape_pct}, TFT's by
\nv{mech.drift.grid.TFT.f0.1.mean_dwape_pct} and TiDE's by
\nv{mech.drift.grid.TiDE.f0.1.mean_dwape_pct}; with the oldest half the figures are
\nv{mech.drift.grid.PLANET.f0.5.mean_dwape_pct}, \nv{mech.drift.grid.CalTST.f0.5.mean_dwape_pct},
\nv{mech.drift.grid.PatchTST.f0.5.mean_dwape_pct},
\nv{mech.drift.grid.NBEATS.f0.5.mean_dwape_pct}, \nv{mech.drift.grid.TFT.f0.5.mean_dwape_pct}
and \nv{mech.drift.grid.TiDE.f0.5.mean_dwape_pct}, and the full-data ordering of
Table~\ref{tab:grids-v2} returns. Under the harshest condition, little and old data, the two
calendar-tier models are the most robust and TiDE the least; with half the history every
method degrades by a comparable amount. The per-window solve follows the level and the
seasonal shape of the current window, but the weather response and the residual network are
frozen, and it is the residual that drifts: the lockdown weeks of spring 2020 are not a
calendar effect. This is the condition that the hypothesis of gradient-free test-time
adaptation was meant to address, and the tier alone does not address it; we report this as the
boundary of the mechanism, not as a failure of the model.

\paragraph{The covariate pathway, isolated} The leading model and the proposed one differ in
two things at once: the architecture and the way the covariate block enters it (TiDE's native
per-step encoder against the shared head). Running TiDE with its native pathway switched off and
the block routed through the same head as PLANET, at TiDE's own selected settings, isolates the
pathway. TiDE via head changes TiDE's WAPE by \nv{path.grid.TiDEh.mean_dwape_vs_TiDE_pct} on
average over the \nv{path.grid.TiDEh.n_cells_vs_TiDE} grid cells (worse in
\nv{path.grid.TiDEh.n_worse_vs_TiDE}, significantly so in \nv{path.grid.TiDEh.n_sig_worse}) and sits \nv{path.grid.TiDEh.mean_dwape_vs_PLANET_pct} behind PLANET on average (worse in
\nv{path.grid.TiDEh.n_worse_vs_PLANET} of \nv{path.grid.TiDEh.n_cells_vs_PLANET} cells): the
part of TiDE's margin that the native pathway accounts for is the first number, and the second
says that once both models read the covariates through the same head, the model with the tier
is ahead.

\subsection{Sensitivity analyses}
\label{sec:sensitivity}
Five checks probe the protocol choices that the referees of the study could reasonably question.
\emph{Correlated weather error.} The headline proxy for a day-ahead weather forecast adds
independent noise at every hour, which is kinder than a real forecast whose errors persist along
the horizon (the daily mean of white noise nearly averages out). Repeating the R2 cells of TiDE
and PLANET on DE, FR and ES with an AR(1) error path (correlation time 12\,h, the same marginal
standard deviation at every lead) raises their WAPE by \nv{exo.ar.grid.TiDE.mean_dwape_pct}
and \nv{exo.ar.grid.PLANET.mean_dwape_pct} on average over \nv{exo.ar.grid.TiDE.n_cells} cells
each, most on France at two weeks (\nv{exo.ar.OPSD-FR.TiDE.H336.dwape_pct} for TiDE,
\nv{exo.ar.OPSD-FR.PLANET.H336.dwape_pct} for PLANET): the model that leans harder on the
weather loses more when its errors persist, and PLANET's deficit against TiDE on these cells
narrows from \nv{exo.ar.grid.gap_white_mean_pct} under white noise to
\nv{exo.ar.grid.gap_PLANET_vs_TiDE_mean_pct} under correlated noise. If the white-noise proxy
of the headline favours either model, it favours the baseline.
\emph{Weather representation.} With the population-weighted mean of six metropolitan points in
place of the single interior point, on the three systems where the weather carries the most
load (FR, ES, IT), both models gain on average, TiDE by \nv{mp.grid.TiDE.mean_dwape_pct} and
PLANET by \nv{mp.grid.PLANET.mean_dwape_pct} over \nv{mp.grid.TiDE.n_cells} cells each (at most
\nv{mp.grid.TiDE.min_dwape_pct} and \nv{mp.grid.PLANET.min_dwape_pct}), so PLANET's deficit
against TiDE on these three systems widens from \nv{mp.grid.gap_single_mean_pct} to
\nv{mp.grid.gap_PLANET_vs_TiDE_mean_pct}: a richer weather input helps the model that reads it
natively more, and the single point of the headline understates the baseline's margin on the
weather-driven systems rather than PLANET's.
\emph{Learning rate per regime.} Re-selecting the learning rate under R0 and R1 on Germany
(Section~\ref{sec:setup}) changes the choice for \nv{lrsel.R0.n_changed} and
\nv{lrsel.R1.n_changed} of \nv{lrsel.R0.n_models} headline models and changes their validation
loss by \nv{lrsel.R0.mean_val_dmse_pct} and \nv{lrsel.R1.mean_val_dmse_pct} on average (at most
\nv{lrsel.R0.min_val_dmse_pct} and \nv{lrsel.R1.min_val_dmse_pct}). Re-running the Germany
cells of the models whose choice changed with their own rate changes their test WAPE by
\nv{lrsel.R0.grid.mean_dwape_own_pct} under R0 and \nv{lrsel.R1.grid.mean_dwape_own_pct} under
R1 (\nv{lrsel.R0.grid.n_better_own} of \nv{lrsel.R0.grid.n_cells_own} and
\nv{lrsel.R1.grid.n_better_own} of \nv{lrsel.R1.grid.n_cells_own} cells improve), and the regime
steps over those models move from \nv{lrsel.r2.OPSD-DE.R1_vs_R0.mean_dwape_pct} and
\nv{lrsel.r2.OPSD-DE.R2_vs_R1.mean_dwape_pct} at the R2 rate to
\nv{lrsel.own.OPSD-DE.R1_vs_R0.mean_dwape_pct} and
\nv{lrsel.own.OPSD-DE.R2_vs_R1.mean_dwape_pct} with the own rates. The reading of
Section~\ref{sec:info} survives: the calendar step shrinks once every model trains at its own
rate but remains several times the weather step on this calendar-driven system, and the
combined change from R0 to R2 over the same cells,
\nv{lrsel.own.OPSD-DE.R2_vs_R0.mean_dwape_pct} against
\nv{lrsel.r2.OPSD-DE.R2_vs_R0.mean_dwape_pct} at the R2 rate, still exceeds the gap between
PLANET and the median learned method within R2 (\nv{grid.median_gap_wape_pct} over the eight
systems).
\emph{Training budget.} Tripling the budget (cap 90 epochs, patience 10) for the headline
models on Germany changes WAPE by \nv{budget.OPSD-DE.mean_dwape_pct} on average over
\nv{budget.OPSD-DE.n_cells} cells (range \nv{budget.OPSD-DE.min_dwape_pct} to
\nv{budget.OPSD-DE.max_dwape_pct}), leaves the per-horizon ranking at a Spearman correlation of
\nv{budget.OPSD-DE.mean_spearman} with the headline ranking and changes the best model in
\nv{budget.OPSD-DE.n_best_changed} of \nv{budget.OPSD-DE.n_horizons} horizons; on France the
corresponding figures are \nv{budget.OPSD-FR.mean_dwape_pct} over \nv{budget.OPSD-FR.n_cells}
cells, \nv{budget.OPSD-FR.mean_spearman} and \nv{budget.OPSD-FR.n_best_changed} of
\nv{budget.OPSD-FR.n_horizons}, the change being at one day, where
\nv{budget.OPSD-FR.H24.best_large} overtakes \nv{budget.OPSD-FR.H24.best_small}.
\emph{Forecast combination.} Whether the two best models are complementary is a question a
benchmark should answer. At the seed whose test forecasts are stored (\nv{comb.seed}), the
equal-weight mean of PLANET and TiDE has a lower WAPE than TiDE alone in
\nv{comb.grid.Comb2.n_better_TiDE} of \nv{comb.grid.n_cells} cells, significantly so in
\nv{comb.grid.Comb2.n_sig_better_TiDE} (Diebold--Mariano with the bootstrap standard error),
by \nv{comb.grid.Comb2.mean_dwape_vs_TiDE_pct} on average, and than PLANET alone in
\nv{comb.grid.Comb2.n_better_PLANET}. The mean of the three models with the lowest validation
loss in each cell, chosen without the test data, changes TiDE's WAPE by
\nv{comb.grid.CombTop3.mean_dwape_vs_TiDE_pct} and includes PLANET in
\nv{comb.grid.CombTop3.n_cells_with_PLANET} cells; the mean of all ten headline models changes
it by \nv{comb.grid.CombAll.mean_dwape_vs_TiDE_pct}. The combinations are not entered in the
headline table, since a single-seed combination cannot be compared with the three-seed means
of its members; the day-ahead table reports the combination of the seed ensembles under the
operator protocol (Section~\ref{sec:dayahead}).

\subsection{Probabilistic forecasts}
\label{sec:prob-results}
\tabv{tab_prob_v2}
\tabv{tab_prob_v2_b}
\begin{figure}[!t]
\centering
\figv{fig_prob_v2}{\columnwidth}
\caption{Day-ahead probabilistic forecasts under R2: CRPS of the seven-level quantile head
(bars; calibration adjusts only the interval bounds) and empirical coverage of the nominal
90\,\% interval before calibration, after CQR and after adaptive conformal inference
(markers), per system.}
\label{fig:prob}
\end{figure}

Every headline model with a quantile head is calibrated by the same two-stage procedure, so
Table~\ref{tab:prob-v2} compares forecasters rather than calibrators. Adaptive conformal
inference brings every method to nominal coverage (PICP after ACI between
\nv{prob.grid.min_picp_aci} and \nv{prob.grid.max_picp_aci} across methods and systems); what
differs is sharpness at that coverage, and it orders the methods as the point grid does. TiDE
has the lowest CRPS on \nv{prob.grid.TiDE.n_best_crps} of \nv{prob.grid.n_systems} systems and
PLANET on \nv{prob.grid.PLANET.n_best_crps}, Germany and France. The mean interval width after
calibration is \nv{prob.grid.TiDE.mean_mpiw_aci}\,MW for TiDE,
\nv{prob.grid.PLANET.mean_mpiw_aci} for PLANET, \nv{prob.grid.NBEATS.mean_mpiw_aci} for
N-BEATS and \nv{prob.grid.CalTST.mean_mpiw_aci} for PatchTST+Cal, and the Winkler scores order
them the same way (\nv{prob.grid.TiDE.mean_winkler_aci}, \nv{prob.grid.PLANET.mean_winkler_aci},
\nv{prob.grid.NBEATS.mean_winkler_aci}, \nv{prob.grid.CalTST.mean_winkler_aci}). TFT, the one
model of the set designed around quantile output, ends with DLinear at the highest CRPS
(\nv{prob.grid.TFT.mean_crps_cqr} and \nv{prob.grid.DLinear.mean_crps_cqr}\,MW) and, with
PatchTST, the widest calibrated intervals (\nv{prob.grid.TFT.mean_mpiw_aci} and
\nv{prob.grid.PatchTST.mean_mpiw_aci}\,MW): its cost buys nothing here either. The raw
quantile heads are under-covered before calibration for every method
(\nv{prob.grid.PLANET.mean_picp_raw} for PLANET), a reminder that pinball-trained quantiles on
a drifting series are not intervals until they are conformalised, and that the calibrator is
what makes the comparison fair.

\subsection{Zero-shot foundation models}
\label{sec:fm}
\tabv{tab_fm_v2}
\begin{figure}[!t]
\centering
\figv{fig_fm_v2}{\columnwidth}
\caption{Zero-shot foundation models against the best learned model of each cell, WAPE
relative to that model. Both foundation models may have seen the public series in
pre-training; the comparison is a reference, not a contest.}
\label{fig:fm}
\end{figure}

Table~\ref{tab:fm-v2} places the two foundation models beside the best learned model of every
cell under R2, with the covariates supplied through Chronos-2's native past- and
future-covariate interface (run in half precision) and, for TimesFM-2.0, through a NumPy port
of the library's in-context regression on covariates. With the protocol lookback of 336 hours
Chronos-2 is more accurate than the best learned model in
\nv{fm.grid.Chronos2.H24.n_beats_best} of \nv{fm.grid.Chronos2.H24.n_cells} systems at one day
and its WAPE is \nv{fm.grid.Chronos2.H24.mean_dwape_vs_best_pct} relative to it on average;
given a context of 1024 hours, longer than any trained model was allowed, it is ahead in
\nv{fm.grid.L1024.Chronos2.H24.n_beats_best} systems and within
\nv{fm.grid.L1024.Chronos2.H24.mean_dwape_vs_best_pct} on average. The gap then opens with the
horizon, \nv{fm.grid.L1024.Chronos2.H96.mean_dwape_vs_best_pct} at four days and
\nv{fm.grid.L1024.Chronos2.H336.mean_dwape_vs_best_pct} at two weeks with the long context
(\nv{fm.grid.Chronos2.H336.mean_dwape_vs_best_pct} with the short), exactly as it does for the
learned calendar maps, which is consistent with a foundation model being a very well trained
window-relative forecaster. TimesFM-2.0 is behind the best learned model in every cell
(\nv{fm.grid.TimesFM2.H24.mean_dwape_vs_best_pct} at one day and
\nv{fm.grid.TimesFM2.H336.mean_dwape_vs_best_pct} at two weeks with the protocol lookback), and
its R2 error exceeds its own R0 error in every cell of Table~\ref{tab:fm-v2}, which points at
the NumPy covariate port rather than at the model; its R0 rows are the ones to read.
Because the pre-training corpora of both models include public electricity series, any
advantage they show is an upper bound on what a clean zero-shot model would achieve, and any
deficit is a lower bound.

\subsection{Multivariate benchmarks}
\label{sec:bench}
\tabv{tab_bench_v2}
On the Electricity panel, scored as the literature scores it (z-scored MSE, all 321 clients,
regime R1 since the panel has no weather), PLANET's mean MSE rank is
\nv{R1.ECL.PLANET.meanrank} over \nv{R1.ECL.n_cells} horizon cells of
\nv{R1.ECL.n_models} methods, against \nv{R1.ECL.PatchTST.meanrank} for PatchTST,
\nv{R1.ECL.CalTST.meanrank} for PatchTST+Cal, \nv{R1.ECL.DLinear.meanrank} for DLinear,
\nv{R1.ECL.iTransformer.meanrank} for iTransformer, \nv{R1.ECL.TiDE.meanrank} for TiDE and
\nv{R1.ECL.NBEATS.meanrank} for N-BEATS, and it is retained by the confidence set in
\nv{R1.mcs.ECL.PLANET.n_in} of those cells; the compact models are within a few percent of one
another in MSE at four days, and the calendar tier is neither help nor harm on this panel. On
the four ETT transformer-station series, which carry an industrial load with weak calendar
structure and a Chinese holiday calendar, the tier is expected to add little and does not: the
best mean ranks over the \nv{R1.ett.n_cells} station--horizon cells belong to the linear RLinear
(\nv{R1.ett.RLinear.meanrank}) and to N-BEATS (\nv{R1.ett.NBEATS.meanrank}), and PLANET's
(\nv{R1.ett.PLANET.meanrank}) sits with PatchTST (\nv{R1.ett.PatchTST.meanrank}) and the
frequency-domain FITS (\nv{R1.ett.FITS.meanrank}), ahead of TiDE (\nv{R1.ett.TiDE.meanrank}).
We report the panel because the field expects it and the stations because a calendar prior
that helped everywhere would be suspicious.

\subsection{Cost}
\label{sec:cost}
\tabv{tab_cost_v2}
Table~\ref{tab:cost-v2} reports parameters, training and inference time on the German
system at four days. The tier adds \nv{cost.tier.params_k}k parameters (the per-channel
penalties, the weight network and the state projection) and a Cholesky solve of a $K\times K$
system per window and channel; the \nv{proto.head.params_k}k-parameter head is attached to
every model. Measured against the standalone N-BEATS row of the table, which is wider than the backbone
inside PLANET (PLANET has \nv{cost.PLANET_vs_backbone.params_pct} parameters relative to it),
the overhead is \nv{cost.PLANET_vs_backbone.train_pct} in training time and
\nv{cost.PLANET_vs_backbone.infer_pct} in inference; the tier's own parameter count above is
from the ablation pair at equal width. PLANET trains in
\nv{OPSD-DE.PLANET.H96.train_s}\,s against \nv{OPSD-DE.TiDE.H96.train_s}\,s for TiDE,
\nv{OPSD-DE.TFT.H96.train_s}\,s for TFT and \nv{OPSD-DE.TimeXer.H96.train_s}\,s for TimeXer: the
most accurate model of the benchmark is also among the cheapest, and the proposed model sits
with it rather than with the expensive Transformers. Under the load-only scoring every
channel-independent model, PLANET included, is trained on the load channel alone, which
yields the identical forecast at $1/C$ of the cost of training all channels; the multivariate
models keep all channels because their forecasts depend on them.
